% REPLACES: main.tex lines 828-841 ("The reversal extends beyond" to "also matters.")
% PRESERVES: the figure at lines 843-856 and the atomless-cost paragraph at lines 862-871.
% RETAINS: all four numerical values below from the existing manuscript.
% NEW SYMBOLS: L, the score bound; mathcal C_L, the class of noise densities.
% OPTIONAL NUMBERS: entry_hypsec_r1, entry_t3_r1; the optional sentence is commented out.
% REFERENCES: existing result-proposition-a-5; no new bibliography record.

The reversal holds for a class of noise laws.
Let \(\mathcal C_L\) denote the positive, absolutely continuous densities
with \(|(\log f)'|\le L\) almost everywhere, where \(0<L<1\).
Set \(b=1/L>1\).
The score bound yields the posterior bounds and the global order bound.
\hyperref[result-proposition-a-5]{Proposition A.5} extends the equilibrium
comparison to this class.
Its upper cost restriction uses the supremum of the posterior under full orders.
That supremum need not be attained.

Laplace noise is the extreme member of the class.
Among laws with median zero, it has the thinnest tails at a fixed score bound.
The Laplace and logistic laws used here are symmetric about zero.
Laplace gives the most informative flow experiment at every pure order profile.
The price construction in Online Appendix A.6 preserves this comparison under its stated conditions.
Online Appendix A.6 proves both claims.
\par
Under full orders, its log odds relative to prior odds equal
\((2/b)\max\{-1,\min\{x,1\}\}\).
The expression clips the Gaussian log odds at variance \(b\).
Gaussian noise has no finite score bound.
Its posterior bounds do not lie strictly inside \((0,1)\).

Logistic noise belongs to the class.
At a common scale \(b\), logistic noise equals Laplace noise plus independent noise.
Under the same price construction, the Laplace price experiment therefore Blackwell-dominates the logistic experiment at the same orders.
The logistic posterior approaches its bounds only as flow tends to infinity.
\par
In the strong benchmark economy, entry is 0.301509 under logistic noise,
compared with 0.522757 under Laplace noise.
The logistic flow threshold is 5.424598398, or 1.495369 noise standard deviations from the center.
The same posterior bounds can thus support different participation rates.
The probability of favorable information near the upper bound matters as well.
% OPTIONAL: include the next sentence only after both quantities appear in reconciliation.csv.
% At the same score bound, entry is \NUM{entry_hypsec_r1} under hyperbolic-secant noise and \NUM{entry_t3_r1} under Student \(t_3\) noise.

% REPLACES: main.tex lines 858-860 ("Figure~\ref{fig:3} compares" to "not a general ordering of informativeness.")
Figure~\ref{fig:3} compares the probability that the posterior exceeds
a preparation threshold near \(M\).
Laplace noise assigns positive probability to the upper posterior bound.
Logistic noise assigns none.
The comparison fixes a common scale and traces where posterior mass sits.
The Laplace experiment is more informative.
Informativeness does not order entry.
